\section{Closed-FLRW metric tangent and the canonical $n=2$ scalar/topographic candidate}
\label{sec:closed-flrw-metric-tangent}

A later BHSM audit identified a normalization conflict in an older critical-
fold construction that attempted to use a branch coordinate $X$ as though it
were a local curvature perturbation.  That issue must not be imported into
the cosmological bridge.

\subsection{Do not identify the branch coordinate with a local curvature field}

In the retained BHSM lineage, the quantity
\begin{equation}
X=H^2+a^{-2}
\end{equation}
is an on-shell branch coordinate.  It is not a stored independent local
scalar field $X(x)$.  Moreover, the same value of $X$ occurs in different
geometric branches with different scalar-curvature normalization.  Therefore
an equation of the form
\begin{equation}
\delta R_4=\delta X
\end{equation}
is not a branch-independent metric-tangent law.

The cosmological bridge instead uses the direct Fr\'echet derivative of the
four-dimensional scalar curvature with respect to the metric:
\begin{equation}
\boxed{
D R_g[k]
=
\nabla_\mu\nabla_\nu k^{\mu\nu}
-
\Box\,\mathrm{tr}_g k
-
R_{\mu\nu}k^{\mu\nu}.
}
\label{eq:DR-metric}
\end{equation}
After the scalar gauge quotient, Eq.~\eqref{eq:DR-metric} is projected
directly onto the physical closed-FLRW $n=2$ scalar metric sector.  No
intermediate $X\to R_4$ normalization is introduced.

\paragraph{Theorem 9 (cosmological metric-tangent normalization).}
Let $k_{\mu\nu}^{(2)}$ be a physical $n=2$ scalar metric perturbation of the
closed-FLRW background.  Then the action-level curvature variation entering
the cosmological metric Hessian is uniquely normalized by
\begin{equation}
\delta R_4^{(2)}
=
D R_{\bar g}[k^{(2)}],
\end{equation}
modulo the already-quotiented scalar gauge directions.  Any auxiliary branch
coordinate may be used only after its map to
$D R_{\bar g}[k^{(2)}]$ has been independently derived on that same branch.

\paragraph{Proof.}
Equation~\eqref{eq:DR-metric} is the geometric Fr\'echet derivative of scalar
curvature.  It depends only on the background metric and its tangent
$k_{\mu\nu}^{(2)}$.  A branch coordinate has no independent status in this
variation unless its derivative with respect to the metric is specified.
Therefore the direct metric derivative fixes the normalization without a
branch-dependent scalar intermediary. \hfill$\square$

\subsection{Lowest Robin collar candidate}

The v5.7 scalar/topographic boundary problem selects the lowest admissible
self-adjoint homogeneous collar mode with Robin zero-flux data.  The
corresponding constant collar profile has
\begin{equation}
\lambda_\rho=0.
\end{equation}
For the nonhomogeneous physical $n=2$ extension, the minimal
harmonic-preserving candidate is to retain this lowest radial profile and
extend the $S^3$ harmonic without changing its normalized boundary amplitude.
On that candidate,
\begin{equation}
I_2=1.
\end{equation}

These two values are \emph{candidate domain data}, not a theorem that the full
retained-action trace/extension map must take this form.

Substitution into Eq.~\eqref{eq:C2-ST} gives
\begin{equation}
\boxed{
\mathcal C_{2,\rm ST}^{\rm cand}
=
\frac1{L^2}
\begin{pmatrix}
5 & -1\\
-1 & 13
\end{pmatrix}.
}
\label{eq:C2-ST-candidate}
\end{equation}
Its determinant is
\begin{equation}
\det(L^2\mathcal C_{2,\rm ST}^{\rm cand})=64
\end{equation}
and its eigenvalues are
\begin{equation}
\boxed{
\lambda_\pm
=
\frac{9\pm\sqrt{17}}{L^2},
}
\label{eq:C2-ST-candidate-eigs}
\end{equation}
both strictly positive.

\subsection{What this closes and what remains}

Equation~\eqref{eq:C2-ST-candidate} closes the scalar/topographic side of the
mixed-Hessian calculation on one explicit, self-adjoint, lowest-radial
candidate extension.  It does not prove that the retained action selects that
extension.

The remaining cosmological action object is therefore reduced to
\begin{equation}
\boxed{
\left\{
\mathcal A_g^{(2)},
\mathcal B_{g2},
\mathcal D_{g,\rm phys}^{(2)}
\right\}
}
\label{eq:remaining-metric-data}
\end{equation}
together with action-level confirmation or rejection of the candidate
$(\lambda_\rho,I_2)=(0,1)$.

The next derivation must evaluate
$\mathcal A_g^{(2)}$ and $\mathcal B_{g2}$ directly from the retained action on
the closed-FLRW physical metric tangent.  It must not reuse the invalid
coefficient-one critical-fold map between $X$ and $R_4$.
